\subsection{有界输入 batch 与两遍构建}

若将全库原始和旋转数组都放在 GPU 上，需要 $4Nd+4Nd=8Nd$ bytes。有界暂存限制这些数组，而压缩索引仍需常驻。对 $n_b$ 个流水缓冲、配置的 add 行数 $a_{\rm add}$ 和每数组预算 $S_{\rm stage}$，暂存行数及原始／旋转设备存储为
\begin{equation}
 a_{\rm in}=\max\!\left\{1,\min\!\left(a_{\rm add},
 \left\lfloor\frac{S_{\rm stage}}{4d}\right\rfloor\right)\right\},
 \qquad S_{\rm float}=8n_ba_{\rm in}d.
 \label{eq:4-staging-rows}
\end{equation}
\[
 4d\leq S_{\rm stage},\quad n_b=3,\ S_{\rm stage}=128\,\mathrm{MiB}
 \Longrightarrow S_{\rm float}\leq2n_bS_{\rm stage}=768\,\mathrm{MiB}.
\]
上界只覆盖这些浮点数组，分配和编码器临时区另计。一遍构建按输入顺序编码，再 gather 为列表顺序；两遍构建先分配和排序，再重读并直接编码到最终位置，以额外遍历换取较少的压缩副本。

\par\medskip\noindent\begin{minipage}{\linewidth}
\captionsetup{hypcap=false}
\centering
\begin{tikzpicture}[font=\footnotesize,
 stage/.style={draw=ThesisNavy!55,fill=ThesisNavy!4,rounded corners=2pt,
 text width=2.5cm,align=center,minimum height=11mm,inner sep=3pt},
 arrow/.style={-{Latex[length=2mm]},draw=ThesisNavy!70}]
\node[stage] (a) at (0,0) {Read + rotate};
\node[stage] (b) at (3.15,0) {Assign + encode\\Input-order arrays};
\node[stage] (c) at (6.3,0) {Gather\\Both copies live};
\node[stage] (d) at (9.45,0) {Repack\\Primary words};
\node[stage] (e) at (0,-1.8) {Assign raw rows\\+ sort};
\node[stage] (f) at (3.15,-1.8) {Reread + rotate\\+ encode};
\node[stage] (g) at (6.3,-1.8) {Scatter directly\\List-order arrays};
\node[stage] (h) at (9.45,-1.8) {Repack\\Primary words};
\foreach \x/\y in {a/b,b/c,c/d,e/f,f/g,g/h}\draw[arrow] (\x)--(\y);
\node[anchor=south west,font=\footnotesize\bfseries] at (-1.3,.7) {One pass};
\node[anchor=south west,font=\footnotesize\bfseries] at (-1.3,-1.1) {Two passes};
\end{tikzpicture}

\captionof{figure}[一遍与两遍构建]{一遍与两遍构建.}
\label{fig:2-add-paths}
\end{minipage}\par\medskip

对 Stella，解析模型给出的行数组占用为：一遍 gather 需要 40.59\,GiB，两遍 repack 需要 21.32\,GiB（\cref{tab:4-phase-storage}）。一遍路径同时保留输入顺序与列表顺序的编码；两遍路径直接写入列表顺序。在首遍分配中，质心通过 $\Pi^\top$ 变回原坐标系，因而无需先旋转数据库行。

\subsection{构建各阶段的显存占用}

实现根据行数组及其他分配的预留量，估计一遍构建所需显存：
\begin{equation}
 H_{\rm single}=N(2b_r+3M+20)+2^{30}\quad\text{bytes}.
 \label{eq:4-peak-heuristic}
\end{equation}
十个存活数组贡献
\[
 3NM+2Nb_r+20N=N(3M+2b_r+20).
\]
1\,GiB 预留量是启发式。两遍构建的最终重打包阶段计入 ID、残差、校正项和两份主码，所需存储为
\begin{equation}
 S_{\rm repack,2}=N(b_r+2M+8).
 \label{eq:4-two-pass-repack}
\end{equation}
整个构建峰值还包括模型、排序、上下文、更早的流水／逆置换临时区和后续查询工作区。\Cref{tab:4-phase-storage} 的阶段值单独计算行数组，并非整个进程峰值。

\begin{table}[!htbp]
\centering
\caption[各构建阶段的行数组占用模型]{行数组占用模型（GiB），$d=1{,}024$、$M=128$、$B_r=8$。BGE 使用一遍，Stella 使用两遍。Stella 的一遍列是反事实容量分析。阶段值不含流水、模型、offsets 和运行时成本。}
\label{tab:4-phase-storage}
\small
\begin{tabular}{@{}lrrr@{}}
\toprule
数据集 & 最终行载荷 & 一遍 Gather & 两遍重打包 \\
\midrule
BGE & 10.90 & 23.05 & 12.11 \\
Stella & 19.20 & 40.59 & 21.32 \\
\bottomrule
\end{tabular}
\end{table}

分配乘积缓冲使用独立容量 $a_{\rm assign}$，由 \texttt{JHQ\_ASSIGN\_BATCH} 或源码默认值选择：至少 8,192 个列表时为 32,768，否则为 8,192。其存储为
\begin{equation}
 S_{\rm dots}=4n_b n_{\rm list}a_{\rm assign}\quad\text{bytes}.
 \label{eq:4-assignment-scratch}
\end{equation}
分配容量独立于当前活跃／暂存行数。在 $d=1536/3072$ 下，两个默认值均分配 3\,GiB。BGE／Stella 的 $3\times4=12$\,GiB 分配临时区在 gather 前释放。对 BGE，12\,GiB 临时区与 23.05\,GiB 行数组 gather 不会同时占用显存。峰值是同一时刻存活分配之和的最大值：
\[
 S_{\rm peak}=\max_t\sum_{A\ {\rm live\ at}\ t}|A|.
\]



Stella 冷构建日志选择两遍：\cref{eq:4-peak-heuristic} 的一遍估计为 41.6\,GiB，同一日志报告尾部阶段可用量为 30.6\,GiB。单独采样的常驻进程总量包含上下文和查询工作区：BGE 为 11.93\,GiB，Stella 为 20.44\,GiB（\cref{sec:4-4}）。补充的 20\,ms 设备采样在 BGE add 开头附近观察到 26.08\,GiB 峰值，在 Stella 重打包时观察到 22.58\,GiB 峰值。较小的两遍重打包量与 Stella 较低的实测峰值一致。这些采样设备总量包含上下文，可能遗漏更短暂的峰值。
